\ifdefined\gkver\else\def\gkver{student}\fi
\documentclass[\gkver]{gaokaozhenti}
\newcommand{\ccnum}[1]{{\fontspec{Noto Sans CJK JP}#1}}
\begin{document}

\chapter{2021年福建}

\begin{enumerate}
\item
%% number: 1
%% paperName: 2021年福建高考第 1 题 4 分
%% typeId: 1
%% type: 单选题
%% chapter: 直线运动
%% point: 平均速度和瞬时速度
%% method: 无
%% score: 4
%% degree: 850
%% duplicateId: 0
%% body:
一游客在武夷山九曲溪乘竹筏漂流，途经双乳峰附近的 M 点和玉女峰附近的 N 点，如图所示。已知该游客从 M 点漂流到 N 点的路程为 $5.4 \Ukm$，用时 $1 \Uh$，M、N 间的直线距离为 $1.8 \Ukm$，则从 M 点漂流到 N 点的过程中 \xzanswer{D}

\onepicture[w=7cm]{figs/01.png}

\fourchoices[answer=D]
{该游客的位移大小为 $5.4 \Ukm$}
{该游客的平均速率为 $5.4 \Ums$}
{该游客的平均速度大小为 $0.5 \Ums$}
{若以所乘竹筏为参考系，玉女峰的平均速度为 0}

%% answer: D
%% memo:
\memoanswer{
	A．位移指的是从 M 点漂流到 N 点的有向线段，故位移大小为 $1.8 \Ukm$，故 A 错误；
	
	B．从 M 点漂流到 N 点的路程为 $5.4 \Ukm$，用时 $1 \Uh$，则平均速率为 $\overline{v}_{\text{率}}=\frac{s}{t}=\frac{5.4}{1} \Ukmh$，故 B 错误；
	
	C．该游客的平均速度大小为 $\overline{v}=\frac{x}{t}=\frac{1.8}{1} \Ukmh=0.5 \Ums$，故 C 正确；
	
	D．以玉女峰为参考系，所乘竹筏的平均速度为 $0.5 \Ums$，若以所乘竹筏为参考系，玉女峰的平均速度也为 $0.5 \Ums$，故 D 错误；
	
	故选 C。
}

\item
%% number: 2
%% paperName: 2021年福建高考第 2 题 4 分
%% typeId: 1
%% type: 单选题
%% chapter: 磁场
%% point: 带电粒子在复合场中的运动
%% method: 无
%% score: 4
%% degree: 750
%% duplicateId: 0
%% body:
一对平行金属板中存在匀强电场和匀强磁场，其中电场的方向与金属板垂直，磁场的方向与金属板平行且垂直纸面向里，如图所示。一质子（\ce{^1_1 H}）以速度 $v_{0}$ 自 O 点沿中轴线射入，恰沿中轴线做匀速直线运动。下列粒子分别自 O 点沿中轴线射入，能够做匀速直线运动的是 \xzanswer{B}（所有粒子均不考虑重力的影响）

\onepicture[w=7cm]{figs/02.png}

\fourchoices[answer=B]
{以速度 $\frac{v_{0}}{2}$ 射入的正电子（\ce{^0_1 e}）}
{以速度 $v_{0}$ 射入的电子（\ce{^0_{-1} e}）}
{以速度 $2v_{0}$ 射入的核（\ce{^2_1 H}）}
{以速度 $4v_{0}$ 射入的 $\alpha$ 粒子（\ce{^4_2 He}）}

%% answer: B
%% memo:
\memoanswer{
	质子（\ce{^1_1 H}）以速度 $v_{0}$ 自 O 点沿中轴线射入，恰沿中轴线做匀速直线运动，将受到向上的洛伦兹力和电场力，满足 $qv_{0}B=qE$，解得 $v_{0}=\frac{E}{B}$，即质子的速度满足速度选择器的条件；
	
	A．以速度 $\frac{v_{0}}{2}$ 射入的正电子（\ce{^0_1 e}），所受的洛伦兹力小于电场力，正电子将向下偏转，故 A 错误；
	
	B．以速度 $v_{0}$ 射入的电子（\ce{^0_{-1} e}），依然满足电场力等于洛伦兹力，而做匀速直线运动，即速度选择器不选择电性而只选择速度，故 B 正确；
	
	C．以速度 $2v_{0}$ 射入的核（\ce{^2_1 H}），以速度 $4v_{0}$ 射入的 $\alpha$ 粒子（\ce{^4_2 He}），其速度都不满足速度选择器的条件 $v_{0}=\frac{E}{B}$，故都不能做匀速直线运动，故 CD 错误；
	
	故选 B。
}

\item
%% number: 3
%% paperName: 2021年福建高考第 3 题 4 分
%% typeId: 1
%% type: 单选题
%% chapter: 交流电
%% point: 变压器
%% method: 无
%% score: 4
%% degree: 650
%% duplicateId: 0
%% body:
某住宅小区变压器给住户供电的电路示意图如图所示，图中 $R$ 为输电线的总电阻。若变压器视为理想变压器，所有电表视为理想电表，不考虑变压器的输入电压随负载变化，则当住户使用的用电器增加时，图中各电表的示数变化情况是 \xzanswer{D}

\onepicture[w=8cm]{figs/03.png}

\fourchoices[answer=D]
{$A_{1}$ 增大，$V_{2}$ 不变，$V_{3}$ 增大}
{$A_{1}$ 增大，$V_{2}$ 减小，$V_{3}$ 增大}
{$A_{2}$ 增大，$V_{2}$ 增大，$V_{3}$ 减小}
{$A_{2}$ 增大，$V_{2}$ 不变，$V_{3}$ 减小}

%% answer: D
%% memo:
\memoanswer{
	不考虑变压器的输入电压随负载变化，即变压器原线圈的输入电压 $U_{1}=U_{V1}$ 不变，根据 $\frac{U_{1}}{U_{2}}=\frac{n_{1}}{n_{2}}$ 可知，变压器副线圈的输出电压 $U_{2}=U_{V2}$ 不变；当住户使用的用电器增加时，即用户的总电阻 $R_{\text{灯}}$ 变小，由 $I_{2}=\frac{U_{2}}{R+R_{\text{灯}}}$ 可知，副线圈的电流 $I_{2}=I_{A2}$ 变大，而由 $U_{V3}=U_{2}-I_{2}R$ 可知 $V_{3}$ 减小；由理想变压器的原理 $U_{1}I_{1}=U_{2}I_{2}$ 可知原线圈的电流 $I_{1}=I_{A1}$ 变大；故综合上述分析可知 $A_{1}$ 增大，$A_{2}$ 增大，$V_{2}$ 不变，$V_{3}$ 减小；
	
	故选 D。
}

\item
%% number: 4
%% paperName: 2021年福建高考第 4 题 4 分
%% typeId: 1
%% type: 单选题
%% chapter: 运动和力
%% point: 动量定理
%% method: 建立模型
%% score: 4
%% degree: 650
%% duplicateId: 0
%% body:
福建属于台风频发地区，各类户外设施建设都要考虑台风影响。已知 10 级台风的风速范围为 $24.5 \Ums \sim 28.4 \Ums$，16 级台风的风速范围为 $51.0 \Ums \sim 56.0 \Ums$。若台风迎面垂直吹向一固定的交通标志牌，则 16 级台风对该交通标志牌的作用力大小约为 10 级台风的 \xzanswer{B}

\fourchoices[answer=B]
{2 倍}
{4 倍}
{8 倍}
{16 倍}

%% answer: B
%% memo:
\memoanswer{
	设空气的密度为 $\rho$，风迎面垂直吹向一固定的交通标志牌的横截面积为 $S$，在时间 $\Delta t$ 的空气质量为 $\Delta m=\rho Sv\cdot\Delta t$，假定台风迎面垂直吹向一固定的交通标志牌的末速度变为零，对风由动量定理有 $-F\cdot\Delta t=0-\Delta mv$，可得 $F=\rho Sv^{2}$，10 级台风的风速 $v_{1}\approx25 \Ums$，16 级台风的风速 $v_{2}\approx50 \Ums$，则有 $\frac{F_{2}}{F_{1}}=\frac{v_{2}^{2}}{v_{1}^{2}}\approx4$；
	
	故选 B。
}

\item
%% number: 5
%% paperName: 2021年福建高考第 5 题 6 分
%% typeId: 2
%% type: 多选题
%% chapter: 机械波
%% point: 机械波的产生
%% method: 无
%% score: 6
%% degree: 850
%% duplicateId: 0
%% body:
以声波作为信息载体的水声通信是水下长距离通信的主要手段。2020 年 11 月 10 日，中国载人潜水器“奋斗者”号创造了 10909 米深潜纪录。此次深潜作业利用了水声通信和电磁通信等多种通信方式进行指令传输或数据交换，如图所示。下列说法正确的是 \xzanswer{BD}

\onepicture[w=10cm]{figs/05.jpg}

\fourchoices[answer=BD]
{“奋斗者”号与“探索一号”通信的信息载体属于横波}
{“奋斗者”号与“沧海”号通信的信息载体属于横波}
{“探索一号”与通信卫星的实时通信可以通过机械波实现}
{“探索一号”与“探索二号”的通信过程也是能量传播的过程}

%% answer: BD
%% memo:
\memoanswer{
	A．由题知，“奋斗者”号与“探索一号”通信是通过水声音通信，由下而上，故信息载体属于纵波，故 A 错误；
	
	B．由题知，“奋斗者”号与“沧海”号通信是通过水声音通信，由左向右，故信息载体属于横波，故 B 正确；
	
	C．因为太空中没有介质，故机械波无法传播，所以“探索一号”与通信卫星的实时通信只能通过电磁通信来实现，故 C 错误；
	
	D．在传递信息的过程也是传递能量的过程，故“探索一号”与“探索二号”的通信过程也是能量传播的过程，故 D 正确。
	
	故选 BD。
}

\item
%% number: 6
%% paperName: 2021年福建高考第 6 题 6 分
%% typeId: 2
%% type: 多选题
%% chapter: 磁场
%% point: 磁场的叠加
%% method: 无
%% score: 6
%% degree: 650
%% duplicateId: 0
%% body:
如图，四条相互平行的细长直导线垂直坐标系 $xOy$ 平面，导线与坐标平面的交点为 a、b、c、d 四点。已知 a、b、c、d 为正方形的四个顶点，正方形中心位于坐标原点 $O$，e 为 cd 的中点且在 $y$ 轴上；四条导线中的电流大小相等，其中过 a 点的导线的电流方向垂直坐标平面向里，其余导线电流方向垂直坐标平面向外。则 \xzanswer{BD}

\onepicture[w=5.5cm]{figs/06.png}

\fourchoices[answer=BD]
{$O$ 点的磁感应强度为 0}
{$O$ 点的磁感应强度方向由 $O$ 指向 c}
{e 点的磁感应强度方向沿 $y$ 轴正方向}
{e 点的磁感应强度方向沿 $y$ 轴负方向}

%% answer: BD
%% memo:
\memoanswer{
	AB．由题知，四条导线中的电流大小相等，且到 $O$ 点的距离相等，故四条导线在 $O$ 点的磁感应强度大小相等，根据右手螺旋定则可知，四条导线中在 $O$ 点产生的磁感应强度方向，如图所示
	
	\onepicture[w=5.5cm]{figs/06_an1.png}
	
	由图可知，$B_{b}$ 与 $B_{c}$ 相互抵消，$B_{a}$ 与 $B_{d}$ 合成，根据平行四边形定则，可知 $O$ 点的磁感应强度方向由 $O$ 指向 c，其大小不为零，故 A 错误，B 正确；
	
	CD．由题知，四条导线中的电流大小相等，a、b 到 e 点的距离相等，故 a、b 在 e 点的磁感应强度大小相等，c、d 到 e 点的距离相等，故 c、d 在 e 点的磁感应强度大小相等，根据右手螺旋定则可知，四条导线中在 e 点产生的磁感应强度方向，如图所示
	
	\onepicture[w=5.5cm]{figs/06_an2.png}
	
	由图可知 $B_{c}$ 与 $B_{d}$ 大小相等，方向相反，互相抵消；而 $B_{b}$ 与 $B_{a}$ 大小相等，方向如图所示，根据平行四边形定则，可知两个磁感应强度的合磁感应强度沿 $y$ 轴负方向，故 C 错误，D 正确。
	
	故选 BD。
}

\item
%% number: 7
%% paperName: 2021年福建高考第 7 题 6 分
%% typeId: 2
%% type: 多选题
%% chapter: 电磁感应
%% point: 双杆问题
%% method: 无
%% score: 6
%% degree: 450
%% duplicateId: 0
%% body:
如图，P、Q 是两根固定在水平面内的光滑平行金属导轨，间距为 $L$，导轨足够长且电阻可忽略不计。图中 EFGH 矩形区域有一方向垂直导轨平面向上、感应强度大小为 $B$ 的匀强磁场。在 $t=t_{1}$ 时刻，两均匀金属棒 a、b 分别从磁场边界 EF、GH 进入磁场，速度大小均为 $v_{0}$；一段时间后，流经 a 棒的电流为 0，此时 $t=t_{2}$，b 棒仍位于磁场区域内。已知金属棒 a、b 相同材料制成，长度均为 $L$，电阻分别为 $R$ 和 $2R$，a 棒的质量为 $m$。在运动过程中两金属棒始终与导轨垂直且接触良好，a、b 棒没有相碰，则 \xzanswer{AD}

\onepicture[w=10cm]{figs/07.png}

\fourchoices[answer=AD]
{$t_{1}$ 时刻 a 棒加速度大小为 $\frac{2B^{2}L^{2}v_{0}}{3mR}$}
{$t_{2}$ 时刻 b 棒的速度为 0}
{$t_{1}\sim t_{2}$ 时间内，通过 a 棒横截面的电荷量是 b 棒的 2 倍}
{$t_{1}\sim t_{2}$ 时间内，a 棒产生的焦耳热为 $\frac{2}{9}mv_{0}^{2}$}

%% answer: AD
%% memo:
\memoanswer{
	A．由题知，a 进入磁场的速度方向向右，b 的速度方向向左，根据右手定则可知，a 产生的感应电流方向是 E 到 F，b 产生的感应电流方向是 H 到 G，即两个感应电流方向相同，所以流过 a、b 的感应电流是两个感应电流之和，则有 $I=\frac{2BLv_{0}}{3R}$。对 a，根据牛顿第二定律有 $BIL=ma$，解得 $a=\frac{2B^{2}L^{2}v_{0}}{3mR}$，故 A 正确；
	
	B．根据左手定则，可知 a 受到的安培力向左，b 受到的安培力向右，由于流过 a、b 的电流一直相等，故两个力大小相等，则 a 与 b 组成的系统动量守恒。由题知，$t_{2}$ 时刻流过 a 的电流为零时，说明 a、b 之间的磁通量不变，即 a、b 在 $t_{2}$ 时刻达到了共同速度，设为 $v$。由题知，金属棒 a、b 相同材料制成，长度均为 $L$，电阻分别为 $R$ 和 $2R$，根据电阻定律有 $R=\rho\frac{L}{s}$，$2R=\rho\frac{L}{s'}$，解得 $s'=\frac{1}{2}s$。已知 a 的质量为 $m$，设 b 的质量为 $m'$，则有 $m=\rho_{\text{密}}sL$，$m'=\rho_{\text{密}}s'L$，联立解得 $m'=\frac{1}{2}m$。取向右为正方向，根据系统动量守恒有 $mv_{0}-\frac{1}{2}mv_{0}=\left(m+\frac{1}{2}m\right)v$，解得 $v=\frac{1}{3}v_{0}$，故 B 错误；
	
	C．在 $t_{1}\sim t_{2}$ 时间内，根据 $q=I\Delta t$，因通过两棒的电流时刻相等，所用时间相同，故通过两棒横截面的电荷量相等，故 C 错误；
	
	D．在 $t_{1}\sim t_{2}$ 时间内，对 a、b 组成的系统，根据能量守恒有 $\frac{1}{2}mv_{0}^{2}+\frac{1}{2}\times\left(\frac{1}{2}m\right)v_{0}^{2}=\frac{1}{2}mv^{2}+\frac{1}{2}\times\left(\frac{1}{2}m\right)v^{2}+Q_{\text{总}}$，解得回路中产生的总热量为 $Q_{\text{总}}=\frac{2}{3}mv_{0}^{2}$。对 a、b，根据焦耳定律有 $Q=I^{2}R\Delta t$，因 a、b 流过的电流一直相等，所用时间相同，故 a、b 产生的热量与电阻成正比，即 $Q_{a}:Q_{b}=1:2$，又 $Q_{a}+Q_{b}=Q_{\text{总}}=\frac{2}{3}mv_{0}^{2}$，解得 a 棒产生的焦耳热为 $Q_{a}=\frac{2}{9}mv_{0}^{2}$，故 D 正确。
	
	故选 AD。
}

\item
%% number: 8
%% paperName: 2021年福建高考第 8 题 6 分
%% typeId: 2
%% type: 多选题
%% chapter: 曲线运动
%% point: 天体运动
%% method: 无
%% score: 6
%% degree: 350
%% duplicateId: 0
%% body:
两位科学家因为在银河系中心发现了一个超大质量的致密天体而获得了 2020 年诺贝尔物理学奖。他们对一颗靠近银河系中心的恒星 $S_{2}$ 的位置变化进行了持续观测，记录到的 $S_{2}$ 的椭圆轨道如图所示。图中 $O$ 为椭圆的一个焦点，椭圆偏心率（离心率）约为 0.87。P、Q 分别为轨道的远银心点和近银心点，Q 与 $O$ 的距离约为 $120 \UAU$（太阳到地球的距离为 $1 \UAU$），$S_{2}$ 的运行周期约为 16 年。假设 $S_{2}$ 的运动轨迹主要受银河系中心致密天体的万有引力影响，根据上述数据及日常的天文知识，可以推出 \xzanswer{BCD}

\onepicture[h=5cm]{figs/08.png}

\fourchoices[answer=BCD]
{$S_{2}$ 与银河系中心致密天体的质量之比}
{银河系中心致密天体与太阳的质量之比}
{$S_{2}$ 在 P 点与 Q 点的速度大小之比}
{$S_{2}$ 在 P 点与 Q 点的加速度大小之比}

%% answer: BCD
%% memo:
\memoanswer{
	A．设椭圆的长轴为 $2a$，两焦点的距离为 $2c$，则偏心率 $0.87=\frac{2c}{2a}=\frac{c}{a}$，且由题知，Q 与 $O$ 的距离约为 $120 \UAU$，即 $a-c=120 \UAU$，由此可得出 $a$ 与 $c$，由于 $S_{2}$ 是围绕致密天体运动，根据万有引力定律，可知无法求出两者的质量之比，故 A 错误；
	
	B．根据开普勒第三定律有 $\frac{a^{3}}{T^{2}}=k$，式中 $k$ 是与中心天体的质量 $M$ 有关，且与 $M$ 成正比；所以，对 $S_{2}$ 是围绕致密天体运动有 $\frac{a^{3}}{T_{S_{2}}^{2}}=k_{\text{致}}\propto M_{\text{致}}$，对地球围绕太阳运动有 $\frac{r_{\text{地}}^{3}}{T_{\text{地}}^{2}}=k_{\text{太}}\propto M_{\text{太}}$，两式相比，可得 $\frac{M_{\text{致}}}{M_{\text{太}}}=\frac{a^{3}T_{\text{地}}^{2}}{r_{\text{地}}^{3}T_{S_{2}}^{2}}$，因 $S_{2}$ 的半长轴 $a$、周期 $T_{S_{2}}$，日地之间的距离 $r_{\text{地}}$，地球围绕太阳运动的周期 $T_{\text{地}}$ 都已知，故由上式，可以求出银河系中心致密天体与太阳的质量之比，故 B 正确；
	
	C．根据开普勒第二定律有 $\frac{1}{2}v_{P}(a+c)=\frac{1}{2}v_{Q}(a-c)$，解得 $\frac{v_{P}}{v_{Q}}=\frac{a-c}{a+c}$，因 $a$、$c$ 已求出，故可以求出 $S_{2}$ 在 P 点与 Q 点的速度大小之比，故 C 正确；
	
	D．$S_{2}$ 不管是在 P 点，还是在 Q 点，都只受致密天体的万有引力作用，根据牛顿第二定律有 $G\frac{Mm}{r^{2}}=ma$，解得 $a=\frac{GM}{r^{2}}$，因 P 点到 $O$ 点的距离为 $a+c$，Q 点到 $O$ 点的距离为 $a-c$，解得 $\frac{a_{P}}{a_{Q}}=\frac{(a-c)^{2}}{(a+c)^{2}}$，因 $a$、$c$ 已求出，故 $S_{2}$ 在 P 点与 Q 点的加速度大小之比可以求出，故 D 正确。
	
	故选 BCD。
}

\item
%% number: 9
%% paperName: 2021年福建高考第 9 题 4 分
%% typeId: 3
%% type: 填空题
%% chapter: 原子和原子核
%% point: 半衰期
%% method: 无
%% score: 4
%% degree: 850
%% duplicateId: 0
%% body:
核污水中常含有氚（\ce{^3_1 H}）等放射性核素，处置不当将严重威胁人类安全。氚 $\beta$ 衰变的半衰期长达 12.5 年，衰变方程为 $\ce{^3_1 H}\rightarrow\ce{^A_Z He}+\ce{^0_{-1} e}+\bar{\nu}_{e}$，其中 $\bar{\nu}_{e}$ 是质量可忽略不计的中性粒子，$Z=$ \tkanswer{2}，$A=$ \tkanswer{3}。若将含有质量为 $m$ 的氚的核污水排入大海，即使经过 50 年，排入海中的氚还剩 \tkanswer{$\frac{1}{16}$} $m$（用分数表示）。

%% answer:
\jdanswer{
	2，3，$\frac{1}{16}$
}

%% memo:
\memoanswer{
	根据电荷数守恒和质量数守恒可得 $Z=2$，$A=3$。
	
	经过 50 年，排入海水中的氚的剩余质量为 $m'=m\left(\frac{1}{2}\right)^{\frac{50}{12.5}}=\frac{1}{16}m$。
}

\item
%% number: 10
%% paperName: 2021年福建高考第 10 题 4 分
%% typeId: 3
%% type: 填空题
%% chapter: 气体性质
%% point: 气体图象
%% method: 无
%% score: 4
%% degree: 750
%% duplicateId: 0
%% body:
如图，一定质量的理想气体由状态 A 变化到状态 B，该过程气体对外 \tkanswer{做正功}（填“做正功”“做负功”或“不做功”），气体的温度 \tkanswer{先升高后降低}（填“升高”“降低”“先升高后降低”“先降低后升高”或“始终不变”）。

\onepicture[w=6cm, align=right]{figs/10.png}

%% answer:
\jdanswer{
	做正功，先升高后降低
}

%% memo:
\memoanswer{
	（1）该过程气体体积增大，对外做正功；
	
	（2）由题图可知，从状态 A 到状态 B，$p$ 与 $V$ 的乘积先增大后减小，根据理想气体状态方程 $\frac{pV}{T}=C$ 可知，气体的温度先升高后降低。
}

\item
%% number: 11
%% paperName: 2021年福建高考第 11 题 6 分
%% typeId: 4
%% type: 实验题
%% chapter: 电路
%% point: 电阻定律
%% method: 无
%% score: 6
%% degree: 650
%% duplicateId: 0
%% body:
某实验小组使用多用电表和螺旋测微器测量一长度为 $80.00 \Ucm$ 电阻丝的电阻率，该电阻丝的电阻值约为 $100\sim200 \UO$，材料未知。实验过程如下：
\begin{enumerate}
	\item
	用螺旋测微器测量该电阻丝的直径，示数如图 \ref{2021福建11a} 所示。该电阻丝的直径为 \tkanswer{1.414} $\Umm[a]$。
	
	\item
	对多用电表进行机械调零。
	
	\item
	将多用电表的选择开关旋至 \tkanswer{$\times10$} 倍率的电阻挡（填“$\times1$”“$\times10$”“$\times100$”或“$\times1$ k”）。
	
	\item
	将黑、红表笔短接，调节欧姆调零旋钮，使指针指在电阻挡零刻度线。
	
	\item
	将黑、红表笔并接在待测电阻丝两端，多用电表的示数如图 \ref{2021福建11b} 所示。该电阻丝的电阻值为 \tkanswer{160} $\UO[a]$。
	
	\item
	测量完成之后，将表笔从插孔拔出，并将选择开关旋到“OFF”位置。
	
	\item
	实验测得的该电阻丝电阻率为 \tkanswer{$3.14\times10^{-4}$} $\UOm[a]$（结果保留 3 位有效数字）。
\end{enumerate}

\twopicture[labelA=2021福建11a, labelB=2021福建11b, wA=3.5cm, wB=8cm, align=right]
{figs/11a.png}
{figs/11b.png}

%% answer:
\jdanswer{
\begin{enumerate}
	\item
	1.414
	
	\item
	$\times10$
	
	\item
	160
	
	\item
	$3.14\times10^{-4}$
\end{enumerate}
}

%% memo:
\memoanswer{
	（1）该电阻丝的直径为 $d=1\Umm+41.4\times0.01\Umm=1.414\Umm$。
	
	（3）使用多用电表欧姆挡测电阻时，为了减小误差，应尽可能使指针偏转至刻度盘中央附近，由于该电阻丝的阻值在 $100\sim200 \UO$，而表盘中央刻度在 $15\sim20$ 左右，所以应选择 $\times10$ 倍率的电阻挡；
	
	（5）$15\sim20$ 之间的分度值为 1，所以该电阻丝的电阻值为 $R=16\times10 \UO=160 \UO$；
	
	（7）根据电阻定律有 $R=\rho\frac{L}{S}=\rho\frac{4L}{\pi d^{2}}$，解得该电阻丝的电阻率为 $\rho=\frac{\pi d^{2}R}{4L}\approx3.14\times10^{-4} \UOm$。
}

\item
%% number: 12
%% paperName: 2021年福建高考第 12 题 6 分
%% typeId: 4
%% type: 实验题
%% chapter: 运动和力
%% point: 变加速直线运动
%% method: 无
%% score: 6
%% degree: 550
%% duplicateId: 0
%% body:
某实验小组利用图 \ref{2021福建12a} 所示的实验装置探究空气阻力与速度的关系，实验过程如下：
\begin{enumerate}
	\item
	首先将未安装薄板的小车置于带有定滑轮的木板上，然后将纸带穿过打点计时器与小车相连。
	
	\item
	用垫块将木板一端垫高，调整垫块位置，平衡小车所受摩擦力及其他阻力。若某次调整过程中打出的纸带如图 \ref{2021福建12b} 所示（纸带上的点由左至右依次打出），则垫块应该 \tkanswer{往右移}（填“往左移”“往右移”或“固定不动”）。
	
	\item
	在细绳一端挂上钩码，另一端通过定滑轮系在小车前端。
	
	\item
	把小车靠近打点计时器，接通电源，将小车由静止释放。小车拖动纸带下滑，打出的纸带一部分如图 \ref{2021福建12c} 所示。已知打点计时器所用交流电的频率为 $50 \UHz$，纸带上标出的每两个相邻计数点之间还有 4 个打出的点未画出。打出 F 点时小车的速度大小为 \tkanswer{0.15} $\Ums[a]$（结果保留 2 位小数）。
	
	\item
	保持小车和钩码的质量不变，在小车上安装一薄板。实验近似得到的某时刻起小车 $v-t$ 图像如图 \ref{2021福建12d} 所示，由图像可知小车加速度大小 \tkanswer{逐渐变小}（填“逐渐变大”“逐渐变小”或“保持不变”）。据此可以得到的实验结论是 \tkanswer{空气阻力随速度增大而增大}。
\end{enumerate}

\fourpicture[labelA=2021福建12a, labelB=2021福建12b, labelC=2021福建12c, labelD=2021福建12d, wA=6.5cm, wB=5cm, wC=6.5cm, wD=4.5cm, align=right]
{figs/12a.png}
{figs/12b.png}
{figs/12c.png}
{figs/12d.png}

%% answer:
\jdanswer{
\begin{enumerate}
	\item
	往右移
	
	\item
	0.15
	
	\item
	逐渐变小，空气阻力随速度增大而增大
\end{enumerate}
}

%% memo:
\memoanswer{
	（2）由题图 \ref{2021福建12b} 可知从左往右点间距逐渐增大，说明小车做加速运动，即平衡摩擦力过度，应减小木板的倾角，即将垫块往右移。
	
	（4）打 F 点时小车的速度大小等于打 E、G 两点之间小车的平均速度大小，即 $v_{F}=\frac{x_{EG}}{2T}=\frac{(6.92-3.85)\times10^{-2}}{2\times0.1} \Ums\approx0.15 \Ums$。
	
	（5）$v-t$ 图像的斜率表示加速度，所以由图像可知小车加速度大小逐渐变小。小车加速度随速度的增大而变小，根据牛顿第二定律可知小车所受合外力 $F$ 随速度的增大而变小。装上薄板后，设小车所受空气阻力大小为 $f$，则 $F=T-f$，而细绳拉力 $T$ 不变，故由此得到的结论是空气阻力随速度增大而增大。
}

\item
%% number: 13
%% paperName: 2021年福建高考第 13 题 12 分
%% typeId: 6
%% type: 计算题
%% chapter: 运动和力
%% point: 牛顿定律的应用
%% method: 无
%% score: 12
%% degree: 550
%% duplicateId: 0
%% body:
一火星探测器着陆火星之前，需经历动力减速、悬停避障两个阶段。在动力减速阶段，探测器速度大小由 $96 \Ums$ 减小到 0，历时 $80 \Us$。在悬停避障阶段，探测器启用最大推力为 $7500 \UN$ 的变推力发动机，在距火星表面约百米高度处悬停，寻找着陆点。已知火星半径约为地球半径的 $\frac{1}{2}$，火星质量约为地球质量的 $\frac{1}{10}$，地球表面重力加速度大小取 $10 \Umsq$，探测器在动力减速阶段的运动视为竖直向下的匀减速运动。求：
\begin{enumerate}
	\item
	在动力减速阶段，探测器的加速度大小和下降距离；
	
	\item
	在悬停避障阶段，能借助该变推力发动机实现悬停的探测器的最大质量。
\end{enumerate}

%% answer:
\jdanswer{
\begin{enumerate}
	\item
	$1.2 \Umsq$，$3840 \Um$
	
	\item
	$1875 \Ukg$
\end{enumerate}
}

%% memo:
\memoanswer{
	（1）设探测器在动力减速阶段所用时间为 $t$，初速度大小为 $v_{1}$，末速度大小为 $v_{2}$，加速度大小为 $a$，由匀变速直线运动速度公式有 $v_{2}=v_{1}-at$ ①，代入题给数据得 $a=1.2 \Umsq$ ②。
	
	设探测器下降的距离为 $s$，由匀变速直线运动位移公式有 $s=v_{1}t-\frac{1}{2}at^{2}$ ③，联立②③式并代入题给数据得 $s=3840 \Um$ ④。
	
	（2）设火星的质量、半径和表面重力加速度大小分别为 $M_{\text{火}}$、$r_{\text{火}}$ 和 $g_{\text{火}}$，地球的质量、半径和表面重力加速度大小分别为 $M_{\text{地}}$、$r_{\text{地}}$ 和 $g_{\text{地}}$，由牛顿运动定律和万有引力定律，对质量为 $m$ 的物体有 $\frac{GM_{\text{火}}m}{r_{\text{火}}^{2}}=mg_{\text{火}}$ ⑤，$\frac{GM_{\text{地}}m}{r_{\text{地}}^{2}}=mg_{\text{地}}$ ⑥，式中 $G$ 为引力常量。设变推力发动机的最大推力为 $F$，能够悬停的火星探测器最大质量为 $m_{max}$，由力的平衡条件有 $F=m_{max}g_{\text{火}}$ ⑦，联立⑤⑥⑦式并代入题给数据得 $m_{max}=1875 \Ukg$ ⑧。
	
	在悬停避障阶段，该变推力发动机能实现悬停的探测器的最大质量约为 $1875 \Ukg$。
}

\item
%% number: 14
%% paperName: 2021年福建高考第 14 题 12 分
%% typeId: 6
%% type: 计算题
%% chapter: 功和能
%% point: 动能定理
%% method: 无
%% score: 12
%% degree: 450
%% duplicateId: 0
%% body:
如图 \ref{2021福建14a}，一倾角 $37\Udu$ 的固定斜面的 AB 段粗糙，BC 段光滑。斜面上一轻质弹簧的一端固定在底端 C 处，弹簧的原长与 BC 长度相同。一小滑块在沿斜面向下的拉力 $T$ 作用下，由 A 处从静止开始下滑，当滑块第一次到达 B 点时撤去 $T$。$T$ 随滑块沿斜面下滑的位移 $s$ 的变化关系如图 \ref{2021福建14b} 所示。已知 AB 段长度为 $2 \Um$，滑块质量为 $2 \Ukg$，滑块与斜面 AB 段的动摩擦因数为 0.5，弹簧始终在弹性限度内，重力加速度大小取 $10 \Umsq$，$\sin37\Udu=0.6$。求：
\begin{enumerate}
	\item
	当拉力为 $10 \UN$ 时，滑块的加速度大小；
	
	\item
	滑块第一次到达 B 点时的动能；
	
	\item
	滑块第一次在 B 点与弹簧脱离后，沿斜面上滑的最大距离。
\end{enumerate}

\twopicture[labelA=2021福建14a, labelB=2021福建14b, wA=6cm, wB=6cm]
{figs/14a.png}
{figs/14b.png}

%% answer:
\jdanswer{
\begin{enumerate}
	\item
	$7 \Umsq$
	
	\item
	$26 \UJ$
	
	\item
	$1.3 \Um$
\end{enumerate}
}

%% memo:
\memoanswer{
	（1）设小滑块的质量为 $m$，斜面倾角为 $\theta$，滑块与斜面间的动摩擦因数为 $\mu$，滑块受斜面的支持力大小为 $N$，滑动摩擦力大小为 $f$，拉力为 $10 \UN$ 时滑块的加速度大小为 $a$。由牛顿第二定律和滑动摩擦力公式有 $T+mg\sin\theta-f=ma$ ①，$N-mg\cos\theta=0$ ②，$f=\mu N$ ③，联立①②③式并代入题给数据得 $a=7 \Umsq$ ④。
	
	（2）设滑块在 AB 段运动的过程中拉力所做的功为 $W$，由功的定义有 $W=T_{1}s_{1}+T_{2}s_{2}$ ⑤，式中 $T_{1}$、$T_{2}$ 和 $s_{1}$、$s_{2}$ 分别对应滑块下滑过程中两阶段所受的拉力及相应的位移大小。依题意，$T_{1}=8 \UN$，$s_{1}=1 \Um$，$T_{2}=10 \UN$，$s_{2}=1 \Um$。设滑块第一次到达 B 点时的动能为 $E_{k}$，由动能定理有 $W+(mg\sin\theta-f)(s_{1}+s_{2})=E_{k}-0$ ⑥，联立②③⑤⑥式并代入题给数据得 $E_{k}=26 \UJ$ ⑦。
	
	（3）由机械能守恒定律可知，滑块第二次到达 B 点时，动能仍为 $E_{k}$。设滑块离 B 点的最大距离为 $s_{max}$，由动能定理有 $-(mg\sin\theta+f)s_{max}=0-E_{k}$ ⑧，联立②③⑦⑧式并代入题给数据得 $s_{max}=1.3 \Um$ ⑨。
}

\item
%% number: 15
%% paperName: 2021年福建高考第 15 题 16 分
%% typeId: 6
%% type: 计算题
%% chapter: 电场
%% point: 带电粒子的直线运动
%% method: 无
%% score: 16
%% degree: 350
%% duplicateId: 0
%% body:
如图 \ref{2021福建15a}，同一竖直平面内 A、B、M、N 四点距 $O$ 点的距离均为 $\sqrt{2}L$，$O$ 为水平连线 AB 的中点，M、N 在 AB 连线的中垂线上。A、B 两点分别固定有一点电荷，电荷量均为 $Q$（$Q>0$）。以 $O$ 为原点，竖直向下为正方向建立 $x$ 轴。若取无穷远处为电势零点，则 ON 上的电势 $\varphi$ 随位置 $x$ 的变化关系如图 \ref{2021福建15b} 所示。一电荷量为 $Q$（$Q>0$）的小球 $S_{1}$ 以一定初动能从 M 点竖直下落，一段时间后经过 N 点，其在 ON 段运动的加速度大小 $a$ 随位置 $x$ 的变化关系如图 \ref{2021福建15c} 所示。图中 $g$ 为重力加速度大小，$k$ 为静电力常量。
\begin{enumerate}
	\item
	求小球 $S_{1}$ 在 M 点所受电场力大小。
	
	\item
	当小球 $S_{1}$ 运动到 N 点时，恰与一沿 $x$ 轴负方向运动的不带电绝缘小球 $S_{2}$ 发生弹性碰撞。已知 $S_{1}$ 与 $S_{2}$ 的质量相等，碰撞前、后 $S_{1}$ 的动能均为 $\frac{4kQ^{2}}{3L}$，碰撞时间极短。求碰撞前 $S_{2}$ 的动量大小。
	
	\item
	现将 $S_{2}$ 固定在 N 点，为保证 $S_{1}$ 能运动到 N 点与之相碰，$S_{1}$ 从 M 点下落时的初动能须满足什么条件？
\end{enumerate}

\threepicture[labelA=2021福建15a, labelB=2021福建15b, labelC=2021福建15c, hA=3.5cm, hB=3.5cm, hC=3.5cm, align=right]
{figs/15a.png}
{figs/15b.png}
{figs/15c.png}

%% answer:
\jdanswer{
\begin{enumerate}
	\item
	$\frac{\sqrt{2}kQ^{2}}{4L^{2}}$
	
	\item
	$\frac{8kQ^{2}\sqrt{gL}}{9gL^{2}}$
	
	\item
	$E_{k}>\frac{(13-8\sqrt{2})kQ^{2}}{27L}$
\end{enumerate}
}

%% memo:
\memoanswer{
	（1）设 A 到 M 点的距离为 $R_{M}$，A 点的电荷对小球 $S_{1}$ 的库仑力大小为 $F_{A}$，由库仑定律有 $F_{A}=\frac{kQ^{2}}{R_{M}^{2}}$ ①。设小球 $S_{1}$ 在 M 点所受电场力大小为 $F_{M}$，由力的合成有 $F_{M}=2F_{A}\sin45\Udu$ ②。联立①②式，由几何关系并代入数据得 $F_{M}=\frac{\sqrt{2}kQ^{2}}{4L^{2}}$ ③。
	
	（2）设 $O$ 点下方 $\frac{L}{2}$ 处为 C 点，A 与 C 的距离为 $R_{C}$，小球 $S_{1}$ 在 C 处所受的库仑力大小为 $F_{C}$，由库仑定律和力的合成有 $F_{C}=2\frac{kQ^{2}}{R_{C}^{2}}\sin\theta$ ④，式中 $\sin\theta=\frac{OC}{R_{C}}$。设小球 $S_{1}$ 的质量为 $m_{1}$，小球 $S_{1}$ 在 C 点的加速度大小为 $a$，由牛顿第二定律有 $F_{C}+m_{1}g=m_{1}a$ ⑤，由图 \ref{2021福建15c} 可知，式中 $a=2g$，联立④⑤式并代入数据得 $m_{1}=\frac{8kQ^{2}}{27gL^{2}}$ ⑥。
	
	设 $S_{2}$ 的质量为 $m_{2}$，碰撞前、后 $S_{1}$ 的速度分别为 $v_{1}$、$v_{1}'$，$S_{2}$ 碰撞前、后的速度分别为 $v_{2}$、$v_{2}'$，取竖直向下为正方向。由动量守恒定律和能量守恒定律有 $m_{1}v_{1}+m_{2}v_{2}=m_{1}v_{1}'+m_{2}v_{2}'$ ⑦，$\frac{1}{2}m_{1}v_{1}^{2}+\frac{1}{2}m_{2}v_{2}^{2}=\frac{1}{2}m_{1}v_{1}'^{2}+\frac{1}{2}m_{2}v_{2}'^{2}$ ⑧。设小球 $S_{2}$ 碰撞前的动量为 $p_{2}$，由动量的定义有 $p_{2}=m_{2}v_{2}$ ⑨。依题意有 $\frac{1}{2}m_{1}v_{1}^{2}=\frac{1}{2}m_{1}v_{1}'^{2}=\frac{4kQ^{2}}{3L}$，$m_{1}=m_{2}$。联立⑥⑦⑧⑨式并代入数据，得 $p_{2}=-\frac{8kQ^{2}\sqrt{gL}}{9gL^{2}}$ ⑩，即碰撞前 $S_{2}$ 的动量大小为 $\frac{8kQ^{2}\sqrt{gL}}{9gL^{2}}$。
	
	（3）设 $O$ 点上方 $\frac{L}{2}$ 处为 D 点。根据图 \ref{2021福建15c} 和对称性可知，$S_{1}$ 在 D 点所受的电场力大小等于小球的重力大小，方向竖直向上，$S_{1}$ 在此处加速度为 0；$S_{1}$ 在 D 点上方做减速运动，在 D 点下方做加速运动，为保证 $S_{1}$ 能运动到 N 点与 $S_{2}$ 相碰，$S_{1}$ 运动到 D 点时的速度必须大于零。
	
	设 M 点与 D 点电势差为 $U_{MD}$，由电势差定义有 $U_{MD}=\varphi_{M}-\varphi_{D}$ \ccnum{⑪}。设小球 $S_{1}$ 初动能为 $E_{k}$，运动到 D 点的动能为 $E_{kD}$，由动能定理有 $m_{1}g(MO-DO)+QU_{MD}=E_{kD}-E_{k}$ \ccnum{⑫}，$E_{kD}>0$ \ccnum{⑬}。由对称性，D 点与 C 点电势相等，M 点与 N 点电势相等，依据图 \ref{2021福建15b} 所给数据，并联立⑥\ccnum{⑪}\ccnum{⑫}\ccnum{⑬}式可得 $E_{k}>\frac{(13-8\sqrt{2})kQ^{2}}{27L}$ \ccnum{⑭}。
}

\end{enumerate}

\end{document}
